\begin{abstract}
  This thesis investigates approximate nearest-neighbor search as a new application domain for the Tenstorrent architecture, with the goal of evaluating whether it can achieve competitive recall while reducing whole-system energy consumption compared with conventional CPU- and GPU-based systems.
  First, the main families of approximate nearest-neighbor algorithms are analyzed to identify the index structure that best matches Tenstorrent's spatial architecture and programming model. Based on this analysis, an IVF-Flat search pipeline is designed and implemented on Tenstorrent hardware. The implementation is then evaluated in terms of queries per second, recall, whole-system energy consumption, and infrastructure cost, and compared with multicore CPU configurations and an NVIDIA A100 GPU system.
  The results show that the GPU achieves the highest throughput across the evaluated recall range, while the Tenstorrent implementation becomes the most energy-efficient measured configuration at the highest evaluated recall point. A simplified hardware-acquisition and electricity-cost model further identifies the operating regions in which Tenstorrent offers a cost advantage over the CPU and GPU baselines. These findings indicate that Tenstorrent is a promising architecture for high-recall ANN workloads, although further performance improvements are currently limited by the maturity and flexibility of its low-level kernel APIs.
\end{abstract}